% ============================================================================
% ARBEITSBLATT DS 7 - Gruppe B: Cultura y Diversión
% ============================================================================

\documentclass[11pt,a4paper]{article}

\usepackage[utf8]{inputenc}
\usepackage[T1]{fontenc}
\usepackage[ngerman]{babel}
\usepackage[left=2cm,right=2cm,top=1.8cm,bottom=1.5cm]{geometry}
\usepackage{helvet}
\usepackage[table]{xcolor}
\usepackage{tabularx}
\usepackage{array}
\usepackage{enumitem}
\usepackage{tcolorbox}
\usepackage{fancyhdr}
\usepackage{lastpage}
\usepackage{amssymb}

\renewcommand{\familydefault}{\sfdefault}

\definecolor{spanisch}{HTML}{c41e3a}
\definecolor{gruppeB}{HTML}{1565c0}
\definecolor{grau}{HTML}{666666}
\definecolor{hellgrau}{HTML}{f5f5f5}

\pagestyle{fancy}
\fancyhf{}
\fancyhead[L]{\small\textcolor{grau}{Spanisch Spätbeginner -- Gruppe B}}
\fancyhead[R]{\small\textcolor{gruppeB}{AB-07-B}}
\fancyfoot[C]{\small\thepage/\pageref{LastPage}}
\renewcommand{\headrulewidth}{0.5pt}

\newcommand{\luecke}[1]{\underline{\hspace{#1}}}
\newcommand{\punktebox}[1]{\hfill\fbox{\small \_\_/#1}}
\newcommand{\es}[1]{\textit{#1}}

\newtcolorbox{merkkasten}[1][]{
    colback=white,
    colframe=gruppeB,
    fonttitle=\bfseries\color{gruppeB},
    title=#1,
    boxrule=1.5pt,
    arc=0pt,
    left=5pt, right=5pt, top=5pt, bottom=5pt
}

\newtcolorbox{infobox}{
    colback=hellgrau,
    colframe=grau,
    boxrule=0.5pt,
    arc=0pt,
    left=5pt, right=5pt, top=3pt, bottom=3pt
}

\newcounter{aufgabe}
\newcommand{\aufgabe}[2]{%
    \stepcounter{aufgabe}%
    \vspace{0.8em}%
    \noindent\textbf{\theaufgabe.} #1 \punktebox{#2}%
    \vspace{0.3em}%
}

\begin{document}

% ============================================================================
% KOPFBEREICH
% ============================================================================
\noindent
\begin{tabularx}{\textwidth}{@{}Xr@{}}
    {\Large\textbf{\textcolor{gruppeB}{Cultura y Diversión}}} & Name: \luecke{5cm} \\[0.3em]
    {\small DS 7 -- Países, personalidades y descripciones} & Datum: \luecke{3cm} \\
\end{tabularx}

\vspace{0.3em}
\hrule
\vspace{0.8em}

% ============================================================================
% MERKKASTEN: Países hispanohablantes
% ============================================================================
\begin{merkkasten}[Kasten 1: El mundo hispanohablante -- Datos importantes]
\vspace{0.2em}

\begin{tabularx}{\textwidth}{@{}|p{3cm}|X|@{}}
\hline
\rowcolor{hellgrau}
\textbf{Información} & \textbf{Datos} \\
\hline
Número de países & \luecke{2cm} países hispanohablantes \\[0.4em]
\hline
Hablantes totales & Aprox. \luecke{2cm} millones de personas \\[0.4em]
\hline
País más grande & \luecke{3cm} (por superficie) \\[0.4em]
\hline
País más poblado & \luecke{3cm} (por habitantes) \\[0.4em]
\hline
Continentes & Europa, \luecke{5cm}, África \\[0.4em]
\hline
\end{tabularx}

\vspace{0.3em}
\end{merkkasten}

\vspace{0.8em}

% ============================================================================
% MERKKASTEN 2: Adjetivos -- Steigerung und Nuancen
% ============================================================================
\begin{merkkasten}[Kasten 2: Adjetivos -- Intensificadores]
\vspace{0.2em}

\textbf{Um Adjektive zu verstärken:}

\begin{tabularx}{\textwidth}{@{}|X|X|X|@{}}
\hline
\rowcolor{hellgrau}
\textbf{Struktur} & \textbf{Bedeutung} & \textbf{Beispiel} \\
\hline
muy + Adjektiv & sehr & Es \es{muy} inteligente. \\[0.4em]
\hline
bastante + Adjektiv & ziemlich & Es \es{bastante} simpático. \\[0.4em]
\hline
un poco + Adjektiv & ein bisschen & Es \es{un poco} perezoso. \\[0.4em]
\hline
nada + Adjektiv & gar nicht & No es \es{nada} egoísta. \\[0.4em]
\hline
\end{tabularx}

\vspace{0.3em}
\end{merkkasten}

\vspace{1em}

% ============================================================================
% AUFGABEN
% ============================================================================

\aufgabe{Países y datos: Ergänze die Tabelle mit Informationen.}{8}

\begin{tabularx}{\textwidth}{@{}|p{2.5cm}|X|X|p{2.5cm}|@{}}
\hline
\rowcolor{hellgrau}
\textbf{País} & \textbf{Capital} & \textbf{Continente} & \textbf{Moneda} \\
\hline
España & \luecke{2.5cm} & Europa & Euro \\[0.4em]
\hline
México & Ciudad de México & \luecke{2.5cm} & \luecke{2cm} \\[0.4em]
\hline
Argentina & \luecke{2.5cm} & Sudamérica & Peso \\[0.4em]
\hline
Colombia & Bogotá & \luecke{2.5cm} & \luecke{2cm} \\[0.4em]
\hline
Cuba & \luecke{2.5cm} & Caribe & \luecke{2cm} \\[0.4em]
\hline
Chile & Santiago & \luecke{2.5cm} & Peso \\[0.4em]
\hline
Perú & \luecke{2.5cm} & Sudamérica & Sol \\[0.4em]
\hline
Venezuela & Caracas & \luecke{2.5cm} & \luecke{2cm} \\[0.4em]
\hline
\end{tabularx}

\vspace{0.8em}

\aufgabe{Adjektive mit Intensificadores: Bilde Sätze nach dem Muster.}{6}

\textbf{Beispiel:} Lionel Messi / famoso / muy $\rightarrow$ \es{Lionel Messi es muy famoso.}

\vspace{0.3em}
\begin{enumerate}[label=\alph*), leftmargin=2em]
    \item Shakira / talentoso / muy\\[0.3em]
    \luecke{12cm}
    \item Rafael Nadal / trabajador / bastante\\[0.3em]
    \luecke{12cm}
    \item Pablo Picasso / creativo / muy\\[0.3em]
    \luecke{12cm}
    \item Frida Kahlo / original / bastante\\[0.3em]
    \luecke{12cm}
    \item Penélope Cruz / guapa / muy\\[0.3em]
    \luecke{12cm}
    \item Gabriel García Márquez / inteligente / muy\\[0.3em]
    \luecke{12cm}
\end{enumerate}

\vspace{0.8em}

\aufgabe{Personenbeschreibung: Lies die Beschreibung und rate die Person.}{6}

\begin{enumerate}[label=\alph*), leftmargin=2em]
    \item Es un hombre de Argentina. Es muy famoso en el mundo del fútbol. Es bastante bajo pero muy rápido. Juega para un equipo en Estados Unidos.\\[0.3em]
    \textbf{¿Quién es?} \luecke{4cm}

    \item Es una mujer de Colombia. Es muy famosa por su música. Baila muy bien y sus caderas son famosas. Es bastante creativa.\\[0.3em]
    \textbf{¿Quién es?} \luecke{4cm}

    \item Es un hombre de España. Es tenista profesional. Es muy trabajador y bastante serio en la cancha. Ha ganado muchos torneos.\\[0.3em]
    \textbf{¿Quién es?} \luecke{4cm}
\end{enumerate}

\vspace{0.8em}

\aufgabe{Escribir: Describe a una persona famosa hispanohablante.}{8}

\begin{infobox}
\textbf{Incluye:} Nombre -- Nacionalidad -- Profesión -- 3 adjetivos con intensificadores -- Por qué es famoso/-a
\end{infobox}

\vspace{0.5em}
\begin{tabularx}{\textwidth}{@{}|X|@{}}
\hline
\\[0.5em]
\luecke{14cm} \\[0.8em]
\luecke{14cm} \\[0.8em]
\luecke{14cm} \\[0.8em]
\luecke{14cm} \\[0.8em]
\luecke{14cm} \\[0.8em]
\luecke{14cm} \\[0.8em]
\hline
\end{tabularx}

\vspace{0.8em}

\aufgabe{Comparación cultural: Vergleiche ein spanischsprachiges Land mit Deutschland.}{6}

\begin{infobox}
\textbf{Hilfe:} Usa: \es{Mientras que en ... / En cambio ... / También ... / A diferencia de ...}
\end{infobox}

\vspace{0.5em}
\textbf{País elegido:} \luecke{5cm}

\vspace{0.3em}
\begin{tabularx}{\textwidth}{@{}|X|@{}}
\hline
\\[0.5em]
\luecke{14cm} \\[0.8em]
\luecke{14cm} \\[0.8em]
\luecke{14cm} \\[0.8em]
\luecke{14cm} \\[0.8em]
\hline
\end{tabularx}

\vspace{1em}
\hrule
\vspace{0.3em}
\begin{flushright}
\textbf{Gesamt: \luecke{1cm} / 34 Punkte}
\end{flushright}

\end{document}
